\emph{Step 1 (the limit is a Clifford algebra, for either parity).} By the definition of the tower (\claimref{obj.tower.tower}), $T_{k} = \Cl(V_{k})$ with $V_{k} = \R^{p,q} \perp H_{1} \perp \cdots \perp H_{k}$, and $\iota_{k}$ is induced by the inclusion $V_{k} \subset V_{k+1}$ (\claimref{std-algebra.functoriality}); each $\iota_{k}$ carries the blade basis of $T_{k}$ into that of $T_{k+1}$ and is injective. Every element of $L_{p,q}$ lies in some floor, so $L_{p,q}$ is generated by its frame $\mathcal{F} = \{x_{i}\} \cup \{e_{j}, f_{j}\}$. We use the universal property of each floor (\claimref{std-algebra.existence-uniqueness}): if $g_{1}, \dots, g_{N}$ is the orthonormal frame of $T_{k}$ and $y_{1}, \dots, y_{N}$ are elements of a unital algebra with $y_{a}^{2} = g_{a}^{2}$ and $y_{a}y_{b} = -y_{b}y_{a}$ for $a \ne b$, there is a unique unital homomorphism $T_{k} \to$ that algebra with $g_{a} \mapsto y_{a}$. The argument of the second paragraph of the proof of Theorem \claimref{tower.algebraic-limit-even-charge}, which does not use the parity of $p - q$, identifies $L_{p,q}$ with $\Cl(\R^{p,q} \perp H_{\infty})$; the steps below need only the universal property of the floors.

\emph{Step 2 (the frame bijection, (i)).} Let $\sigma$ be the bijection of (i) from $\mathcal{F}$ to the frame $\mathcal{F}' = \{e'_{j}, f'_{j}\}$ of $L_{0,0}$: it lists the positive generators $x_{1}, \dots, x_{p}, e_{1}, e_{2}, \dots$ in order onto $e'_{1}, e'_{2}, \dots$ and the negative generators $x_{p+1}, \dots, x_{n}, f_{1}, f_{2}, \dots$ in order onto $f'_{1}, f'_{2}, \dots$. It preserves squares, and distinct frame elements of $\mathcal{F}'$ anticommute, so for each $k$ the images of the frame of $T_{k}$ satisfy the relations of that frame, and the universal property gives $\Phi_{k} \colon T_{k} \to L_{0,0}$. The images lie in $T'_{k + \max(p,q)}$, since $x_{i} \mapsto e'_{i}$ with $i \le p$, $x_{p+i} \mapsto f'_{i}$ with $i \le q$, $e_{j} \mapsto e'_{p+j}$ and $f_{j} \mapsto f'_{q+j}$ with $j \le k$. As $\Phi_{k+1} \circ \iota_{k}$ and $\Phi_{k}$ agree on the frame of $T_{k}$, they are equal, and the $\Phi_{k}$ assemble to $\Phi \colon L_{p,q} \to L_{0,0}$. In the same way $\sigma^{-1}$ gives $\Psi \colon L_{0,0} \to L_{p,q}$, carrying $T'_{k}$ into $T_{k}$. Then $\Psi\Phi$ and $\Phi\Psi$ are unital homomorphisms fixing every frame element, and the limits are generated by their frames, so both are identities: $\Phi$ is an isomorphism. By Theorem \claimref{tower.algebraic-limit-even-charge} at $p = q = 0$, $L_{0,0} \cong \Cl(0,0) \otimes \Mat_{2^{\infty}} = \Mat_{2^{\infty}}$, and $L_{0,0} = \Cl(H_{\infty})$ by Step 1. For $n \ge 1$, $B_{\infty}$ is the span of the blades in plane generators alone, so $x_{1} \notin B_{\infty}$ by linear independence of blades and $B_{\infty} \ne L_{p,q}$; the plane subalgebra of $L_{0,0}$ is all of $L_{0,0}$, and $\Phi$ is bijective, so $\Phi(B_{\infty}) \ne L_{0,0}$. Since the algebras $L_{p,q}$ for all $(p,q)$ are isomorphic to one algebra, the algebra alone determines neither $n$ nor $r$.

\emph{Step 3 ((ii)).} For $r$ even, Theorem \claimref{tower.algebraic-limit-even-charge} gives $L_{p,q} \cong \Cl(p,q) \otimes \Mat_{2^{\infty}}$, and (i) gives $L_{p,q} \cong \Mat_{2^{\infty}}$. The even residues are of types $\Mat_{k}(\R)$ ($r \equiv 0, 2$) and $\Mat_{k}(\HH)$ ($r \equiv 4, 6$), and both are absorbed.

\emph{Step 4 ((iii)).} Let $r$ be odd, so $n$ is odd and $n \ge 1$. If $q \ge 1$, put $u_{i} = x_{i}x_{n}$ for $i < n$, with $x_{n}^{2} = -1$: then $u_{i}^{2} = -x_{i}^{2}x_{n}^{2} = x_{i}^{2}$ and $u_{i}u_{j} = -u_{j}u_{i}$ for $i \ne j$, so the universal property gives a homomorphism $\Cl(p, q-1) \to \Cl^{0}(p,q)$; it is onto, since $u_{i}u_{j} = -x_{n}^{2}\,x_{i}x_{j}$ gives $x_{i}x_{j} = -x_{n}^{2}\,u_{i}u_{j}$ for $i, j < n$ ($x_{n}^{2} = \pm 1$), and $x_{i}x_{n} = u_{i}$, while every even blade is a product of such pairs; and both sides have dimension $2^{n-1}$, so it is an isomorphism. If $q = 0$, then $x_{n}^{2} = +1$, $u_{i}^{2} = -x_{i}^{2} = -1$, and the same argument gives $\Cl^{0}(p,0) \cong \Cl(0, p-1)$. The residues $r + 1$ and $1 - p$ are even. Write $(p', q')$ for the signature obtained. By (i) twice and Theorem \claimref{tower.algebraic-limit-even-charge} for the even residue,
\[
L_{p,q} \;\cong\; \Mat_{2^{\infty}} \;\cong\; L_{p',q'} \;\cong\; \Cl(p', q') \otimes \Mat_{2^{\infty}} \;\cong\; \Cl^{0}(p,q) \otimes \Mat_{2^{\infty}}.
\]
For the centre: $A \otimes \Mat_{2^{\infty}}$ is the increasing union of $A \otimes \Mat_{2^{k}} = \Mat_{2^{k}}(A)$ along $z \mapsto \operatorname{diag}(z, z)$. A central $z$ lies in some $\Mat_{2^{k}}(A)$ and commutes with the real matrix units of $\Mat_{2^{k}}(\R) \subset \Mat_{2^{k}}(A)$, which forces $z = \operatorname{diag}(a, \dots, a) = a \otimes 1$; commuting with $A \otimes 1$ forces $a \in Z(A)$. Conversely $Z(A) \otimes 1$ commutes with the generating subalgebras $A \otimes 1$ and $1 \otimes \Mat_{2^{\infty}}$. So $Z(A \otimes \Mat_{2^{\infty}}) = Z(A) \otimes 1 \cong Z(A)$. For $L_{p,q} \cong \R \otimes \Mat_{2^{\infty}}$ this is $\R$. For $A = \Cl(p,q)$ with $r$ odd, Lemma \claimref{tower.grade-twist-when-it-inner} gives the centre $\R \oplus \R$ or $\C$, of dimension $2$. Isomorphic algebras have isomorphic centres, so $L_{p,q} \not\cong \Cl(p,q) \otimes \Mat_{2^{\infty}}$. On the floors: $T_{k}$ has odd dimension $n + 2k$, so its volume element $\omega_{T_{k}}$ commutes with every vector of $V_{k}$ and is central, and the centre of $T_{k}$ is $\mathrm{span}\{1, \omega_{T_{k}}\}$: the same lemma applied to $T_{k}$, whose residue is again $r$, gives a two-dimensional centre, which contains this two-dimensional span. The generator $e_{k+1}$ is orthogonal to all $n + 2k$ factors of $\omega_{T_{k}}$, so it anticommutes with it, and $a + b\,\omega_{T_{k}}$ commutes with $e_{k+1}$ only for $b = 0$.

\emph{Step 5 ((iv)).} The frame $*$ of Lemma \claimref{tower.structure-tower-c-norms} is the conjugate-linear antiautomorphism with $g^{*} = g$ for frame generators of square $+1$ and $g^{*} = -g$ for those of square $-1$, compatible with the $\iota_{k}$, hence defined on the complexified limits. Since $\sigma$ preserves squares, $\Phi^{\C} \circ *$ and $* \circ \Phi^{\C}$ are conjugate-linear antihomomorphisms agreeing on the frame, so they are equal: $\Phi^{\C}$ is a $*$-isomorphism $L_{p,q}^{\C} \to L_{0,0}^{\C}$. A C*-norm on $L_{p,q}^{\C}$ is carried by $\Phi^{\C}$ to a C*-norm on $L_{0,0}^{\C}$, and $p - q = 0$ is even, so by Lemma \claimref{tower.structure-tower-c-norms} there is exactly one; hence $L_{p,q}^{\C}$ carries exactly one C*-norm, and $\Phi^{\C}$ extends isometrically to the completions. The completion of $L_{0,0}^{\C}$ is $\Cl(0,0)_{\C} \otimes \CAR = \CAR$ by Theorem \claimref{tower.car-limit-even-charge}.

\emph{The closing sentence.} By (i) to (iii) the algebra is the same for every $(p,q)$, so no invariant of it separates the cases. Theorem \claimref{synthesis.plane-commutant-even-ground}(i) identifies the relative commutant of $B_{\infty}$ in $L_{p,q}$ as $\Cl^{0}(p,q)$ for either parity of $r$, so the pair $(L_{p,q}, B_{\infty})$, unlike the algebra, keeps information about the residue.

The check verifies on truncations, in exact arithmetic, that the frame bijection respects squares and anticommutation and the floor bookkeeping, with $\Psi\Phi$ the identity on generators; the presentation of $\Cl^{0}(p,q)$ as $\Cl(p,q-1)$ or $\Cl(0,p-1)$ for odd $n$; and, for odd $r$, that the centre of each floor is spanned by $1$ and the floor's volume element while the elements of a floor commuting with the next floor's generators are the scalars.
